\section{Introduction}
\label{sec:ch3_introduction}

Edge computing and cloud-native architectures are reshaping the computational landscape by bringing computational resources closer to end-users, enabling low-latency services \cite{shang2023online}. This shift enhances application performance but introduces challenges in optimizing resource utilization due to highly dynamic and fluctuating workloads in edge-cloud environments ~\cite{10229064, 8258257}. Efficient resource management is essential to meet \ac{QoS} requirements while minimizing operational costs~\cite{8258257}.

Accurate workload prediction at the container level has become crucial for proactive orchestration and resource management in edge-cloud systems \cite{9068614}. The widespread adoption of microservice architectures has introduced new challenges, as applications are decomposed into numerous containers with complex dependencies and resource consumption patterns \cite{10029931}. This architectural shift demands more granular predictions for effective container-specific elasticity operations. However, traditional prediction models struggle to capture these complex dynamics, as containerized microservices exhibit unpredictable demand fluctuations and frequent migrations across heterogeneous edge-cloud environments \cite{ouyang2023dynamic, benidis2022deep}.

The state-of-the-art in cloud workload prediction can be categorized into five main approaches \cite{kumar2018workload, tang2019large}. Evolutionary Learning methods integrate neural networks with evolutionary optimization techniques \cite{kumar2021self}. Deep Learning approaches form the predominant category, particularly focusing on temporal dependencies in cloud workloads \cite{zhang2018efficient, zhang2017resource}, though they face challenges with long-range predictions and rapid workload changes \cite{hochreiter1998vanishing}. Hybrid Learning combines multiple algorithms to overcome individual limitations \cite{kardani2020adrl, karim2021bhyprec}, while Ensemble Learning uses multiple predictive models for improved accuracy \cite{singh2014ensemble, iqbal2019adaptive}. Finally, Quantum Learning offers potential for handling the complexity of modern cloud environments \cite{singh2021quantum}. However, most of these approaches rely on classical architectures like \acp{RNN}, while the machine learning literature has advanced with newer architectures.

Recent machine learning advancements have introduced architectures like \acp{GNN}~\cite{yi2024fouriergnn}, temporal \ac{CNN} architectures~\cite{donghao2024moderntcn}, and Transformer-based models~\cite{gort2024forecasting}. These models excel in capturing complex temporal patterns but are often computationally intensive due to their large number of parameters, making them unsuitable for resource-constrained edge devices \cite{meuser2024revisiting}. For example, models like \ac{WGAN}~\cite{arbat2022wassersteinadversarialtransformercloud} and Pathformer~\cite{pathformer2024} contain millions of parameters, requiring large computational resources and high energy consumption, which is impractical for many edge environments \cite{kong2022edge}.

To mitigate these limitations, lightweight \ac{MLP} architectures ~\cite{ekambaram2023tsmixer} have been proposed. One of these models, \ac{SparseTSF}~\cite{sparseTSF} is based on a sparse \ac{MLP} architecture and is designed for fast inference with minimal resource consumption, making it suitable for edge deployment. It achieves performance comparable to larger models while significantly reducing model size. However, \ac{SparseTSF} assumes a fixed periodicity in workload data, limiting its adaptability in environments where workload patterns exhibit dynamic periodicities or abrupt changes.

Models assuming fixed periodicities may fail to capture sudden changes or multiple overlapping patterns ~\cite{sparseTSF}, leading to decreased prediction accuracy and suboptimal resource allocation. Lightweight models that adaptively adjust to changing workload data are needed to improve prediction accuracy and resource efficiency \cite{fan2019urbanedge}. For a detailed review of related work on time-series forecasting models, see Section~\ref{sec:ch2_sota_aero}.

In this chapter, we introduce \ac{AERO}, a novel forecasting model designed to address the challenges of dynamic workload prediction in resource-constrained environments. \ac{AERO} is a lightweight model with fewer than 1,000 parameters, making it suitable for deployment on edge devices with limited computational capacity. It incorporates an adaptive period detection mechanism that dynamically identifies dominant periods in workload data, allowing it to adjust to varying periodic patterns. This adaptability enhances prediction accuracy and resource allocation efficiency compared to models assuming fixed periodicities like \ac{SparseTSF}.

To validate the effectiveness of \ac{AERO}, we used the \ac{COSCO} framework~\cite{tuli2021cosco}, which provides a realistic emulation environment for service orchestration. \ac{COSCO} accurately replicates real-world container behavior, including migration patterns, resource allocation, and system dynamics, enabling thorough evaluation of workload predictors without physical infrastructure costs. The framework's creators have extensively validated it against real fog computing environments, demonstrating high correlation between simulated and actual performance metrics \cite{tuli2021cosco}. As a novel contribution, we extend \ac{COSCO} with a prediction interface.

Our key contributions are:
\begin{itemize}
    \item We propose \ac{AERO}, a lightweight adaptive model for edge-cloud workload prediction that effectively handles workloads with multiple or changing periodicities in multivariate data.
    \item We demonstrate that \ac{AERO} achieves comparable or better accuracy than state-of-the-art models while greatly reducing model size and computational overhead, suiting resource-constrained environments.
    \item We extend the \ac{COSCO} simulation framework with a prediction interface.
    \item Through evaluations on real-world datasets and simulations, we demonstrate \ac{AERO}'s ability to improve resource allocation and reduce operational costs in edge-cloud environments.
\end{itemize}

The remainder of the chapter is organized as follows: Section~\ref{sec:ch3_system_model} presents the system model and formulates the workload prediction problem. Section~\ref{sec:ch3_proposed_solution} describes the architecture and implementation of the \ac{AERO} model in detail. Section~\ref{sec:ch3_evaluation} provides a detailed evaluation of \ac{AERO} compared to existing models. Section~\ref{sec:ch3_simulations} presents detailed simulation results using our extended prediction interface. Section~\ref{sec:ch3_live_deployment} validates the framework in a live edge-cloud deployment. Finally, Section~\ref{sec:ch3_conclusion} concludes the chapter and discusses future work.

